\chapter{\texttt{medium.vacuum} --- coherent propagation in vacuum}
\label{ch:vacuum}

\texttt{medium.vacuum} is the simplest propagation medium in the project:
$n_e\equiv0$ everywhere, so the matter term of the Hamiltonian
(\cref{ch:core-common}) never enters and every formula reduces to the
closed-form vacuum evolutor already derived, in general terms, in the
coherent-propagation scheme of \cref{ch:core-sm}. Its role is twofold:

\begin{itemize}
  \item \textbf{Base case and test bed.} With no density profile to model
    and a fully closed-form evolutor, this is the cheapest possible path
    through the entire propagation stack (\pyclass{PMNS} $\to$
    \pyfunc{kinetic\_eigenvalue\_vector} $\to$ evolutor $\to$ probability
    $\to$ flux, \cref{ch:core-common}), with an exact analytic answer to
    check every other piece against. It is, in practice, the first thing
    exercised when validating a new oscillation preset, a new BSM
    extension, or a refactor of the shared infrastructure: if
    \texttt{medium.vacuum} disagrees with the legacy \texttt{peanuts}
    reference (\cref{sec:vacuum-validation}), nothing built on top of it
    can be trusted either.
  \item \textbf{Physical use cases.} Vacuum propagation is the correct
    description whenever matter effects are genuinely negligible along the
    whole path: reactor and other short-baseline experiments, or any
    scenario where the accumulated matter phase $V\cdot L$ is small enough
    to ignore relative to the vacuum kinetic phases. It is also the
    $n_e\to0$ reference limit that the matter-propagation media of the
    following chapters must recover exactly in that limit.
\end{itemize}

\begin{notebox}
\textbf{The four-module pattern.} \texttt{medium.vacuum} fixes a module
layout that every subsequent medium in this part of the document repeats,
adapted to that medium's own needs: an \texttt{evolutor.py} building the
medium-specific $S$ from \pyclass{OscillationParameters} and geometry
inputs; a \texttt{probability.py} converting that evolutor into transition
and final-state probabilities via the shared
\texttt{core.common.probability} conventions (\cref{ch:core-common}); a
\texttt{flux.py} applying \texttt{core.common.flux} normalisations on top;
and, where a legacy NumPy reference implementation exists to check against,
a \texttt{validation.py} comparing the two numerically. Media with
additional physical structure add further modules on top of this
skeleton (\texttt{profile.py} for a density profile,
\texttt{geometry.py} for path geometry, \texttt{exposure\_*.py} for
detector exposure, ...) without changing this core four-piece contract.
\end{notebox}

\section{\texttt{evolutor.py} --- the vacuum evolutor}

\modulehead{medium/vacuum/evolutor.py}{\pyfunc{vacuum\_evolutor} and
\pyfunc{vacuum\_evolved\_state}: the closed-form coherent vacuum operator
and its action on a state.}

\begin{theorybox}
With $H_{\rm mat}=0$, the reduced and flavour Hamiltonians coincide with
the pure kinetic term, and the general coherent-propagation scheme of
\cref{ch:core-sm} is exact and complete on its own (no reduction trick is
needed, since there is no matter term to factor out):
\begin{equation}
  \keyresult{S(L,E) = U\,\diag\!\bigl(e^{-ik_1x},e^{-ik_2x},e^{-ik_3x}\bigr)\,U^\dagger},
  \qquad x = \frac{L\,[\mathrm{m}]}{L_{\rm scale}},
\end{equation}
where $U$ is the full PMNS matrix (\cref{ch:core-common}), $k_i$ the
dimensionless kinetic eigenvalues (\cref{ch:core-common},
\texttt{potential.py} section), and the baseline $L$ (given in km) is
converted to metres and non-dimensionalised by the same
\code{evolution\_scale\_m} used throughout the project. This is the same
formula the sterile extension recovers automatically for $n=4$
(\cref{ch:core-bsm}), simply by passing a $4\times4$ \pyclass{PMNS\_sterile}
object in place of \pyclass{PMNS\_SM}; nothing in this module is
3-flavour-specific.
\end{theorybox}

\begin{implbox}
\begin{itemize}
  \item \pyfunc{vacuum\_evolutor(oscillation, E\_MeV, L\_km, ...)}: builds
    $S$ exactly as above; broadcasts energy, baseline, and mass-splitting
    batch shapes together before the phase multiplication.
  \item \pyfunc{vacuum\_evolved\_state(nustate, oscillation, E\_MeV, L\_km, ...)}:
    applies \pyfunc{core.common.evolutor.apply\_evolutor\_to\_state}
    (\cref{ch:core-common}) to a coherent flavour-basis initial state.
  \item \code{legacy\_precision} is accepted for call-site symmetry with
    every matter-propagation evolutor in later chapters but is a genuine
    no-op here: there is no matter potential for the flag to affect.
\end{itemize}
\end{implbox}

\section{\texttt{probability.py} --- vacuum probabilities}

{\sloppy\modulehead{medium/vacuum/probability.py}{\pyfunc{vacuum\_\allowbreak{}probability\_\allowbreak{}transition}
(the full transition-probability matrix), \pyfunc{vacuum\_\allowbreak{}probability\_\allowbreak{}state}
(final-state probabilities for a given initial state), and
\pyfunc{vacuum\_\allowbreak{}probability\_\allowbreak{}integrated} (spectrum-weighted energy average).}}

\begin{implbox}
\begin{itemize}
  \item \pyfunc{vacuum\_probability\_transition(oscillation, E\_MeV, L\_km, ...)}:
    builds $S$ and returns $P_{\beta\alpha}=|S_{\beta\alpha}|^2$ via
    \pyfunc{core.common.probability.probability\_transition}
    (\cref{ch:core-common}) -- the full $3\times3$ (or $4\times4$) channel
    matrix, no initial state required.
  \item \pyfunc{vacuum\_probability\_state(nustate, oscillation, E\_MeV, L\_km, massbasis, ...)}:
    dispatches to the coherent or incoherent formula of
    \cref{ch:core-common} (\texttt{evolutor.py}/\texttt{probability.py}
    sections) via \pyfunc{probability\_state}, according to
    \code{massbasis}: \code{False} evolves \code{nustate} as coherent
    flavour amplitudes ($P_\beta=|(S\,\psi)_\beta|^2$); \code{True} (the
    default here) treats it as incoherent mass-basis weights
    ($P_\beta=\sum_i|(S\,U)_{\beta i}|^2w_i$). Either way the function
    returns a final flavour-probability vector, never a flux.
  \item \pyfunc{vacuum\_probability\_integrated(nustate, oscillation, E\_MeV, L\_km, spectrum, ...)}:
    builds \pyfunc{vacuum\_probability\_state} on an energy grid and averages
    it with \pyfunc{core.common.probability.probability\_integrated}
    (\cref{ch:core-common}), weighted by an explicit production
    \code{spectrum} -- there is no default weight, since a flat spectrum
    over an arbitrary grid is a modelling choice.
\end{itemize}
\end{implbox}

\section{\texttt{flux.py} --- vacuum flux}

\modulehead{medium/vacuum/flux.py}{\pyfunc{vacuum\_\allowbreak{}flux\_\allowbreak{}state} and
\pyfunc{vacuum\_\allowbreak{}flux\_\allowbreak{}integrated}: apply \texttt{core.common.flux}
normalisation and energy integration to
\pyfunc{vacuum\_\allowbreak{}probability\_\allowbreak{}state}'s output.}

\begin{implbox}
\pyfunc{vacuum\_flux\_state(nustate, oscillation, E\_MeV, L\_km, flux, spectrum, massbasis, ...)}
adds no new physics: it calls \pyfunc{vacuum\_probability\_state} above,
then \pyfunc{core.common.flux.flux\_state} (\cref{ch:core-common}) to apply
the flux normalisation and optional spectral weight. This is the minimal
instance of the flux pattern that every later medium's \texttt{flux.py}
extends. \pyfunc{vacuum\_flux\_integrated(...)} builds
\pyfunc{vacuum\_flux\_state} on an energy grid and integrates it with
\pyfunc{core.common.flux.flux\_integrated}, returning an unnormalised
physical rate rather than the normalised average of
\pyfunc{vacuum\_probability\_integrated}.
\end{implbox}

\section{\texttt{validation.py} --- cross-check against legacy \texttt{peanuts}}
\label{sec:vacuum-validation}

\modulehead{medium/vacuum/validation.py}{Numerical comparison against the
bundled legacy, NumPy-based \texttt{peanuts} package.}

\begin{implbox}
The legacy vacuum implementation is scalar (one energy, one baseline, one
initial state at a time); this module evaluates both implementations on
matching scalar inputs and compares the results directly, establishing the
pattern reused by \texttt{medium.solar.validation} and
\texttt{medium.earth.validation} (\cref{ch:solar,ch:earth}).
\begin{itemize}
  \item \pyfunc{ensure\_legacy\_importable()} / \pyfunc{legacy\_modules()}:
    locate and import the bundled \code{peanuts} package (kept in the
    repository specifically as this validation reference, not as a runtime
    dependency of \texttt{tpeanuts} itself).
  \item \pyfunc{legacy\_pmns\_from\_torch(...)}: rebuilds a legacy
    \code{peanuts.pmns.PMNS} object from a torch \pyclass{PMNS} instance's
    parameters, so both implementations start from identical inputs.
  \item \pyfunc{compare\_vacuum\_probability\_state\_with\_legacy(...)} and
    \pyfunc{compare\_vacuum\_evolved\_\allowbreak{}state\_with\_\allowbreak{}legacy(...)}:
    run both implementations and return their outputs (and/or a residual)
    for the test suite to assert against.
\end{itemize}
\end{implbox}
